\section{Codifica video}
\subsection{Principi di compressione video}
\subsubsection*{Predizione temporale e decoded frame buffer}
I codificatori video, oltre alla compressione spaziale e alla compressione basata su JPEG, sfruttano anche la ridondanza
temporale tra fotogrammi successivi attraverso perdizioni temporali. Le immagini consecutive di un video, infatti, sono
molto simili tra loro e tale sistema di predizione risulta estremamente efficace per ridurre il bitrate.

In pratica il codificatore calcola i parametri del modello e l'errore di predizione dell'immagine corrente in funzione di
una o più immagini precedenti (o successive) e invia al decodificatore tali informazioni. Il decodificatore avrà un buffer
di immagini decodificate (decoded frame buffer) dove potrà memorizzare temporaneamente le immagini decodificate per poterle
usare come riferimento per la predizione dei fotogrammi successivi.

\subsubsection*{Criticità della predizione temporale e GOP}
Con la predizione temporale si forma una catena continua di dipendenze temporali tra i fotogrammi e si verificano due
criticità principali:
\begin{itemize}
	\item propagazione di errori: se un fotogramma viene perso o danneggiato, tutti i fotogrammi successivi che dipendono
	da esso saranno anch'essi danneggiati
	\item mancanza di accesso casuale: per accedere a un fotogramma intermedio, è necessario decodificare tutti i fotogrammi
	precedenti in un processo lungo e inutilmente costoso.
\end{itemize}
Per mitigare tali problemi, i fotogrammi vengono raggruppati in unità chiamate GOP (Group of Pictures) che rendono fotogrammi
di GOP diversi indipendenti tra loro.

\subsubsection*{Parametri e metriche di prestazione dei codificatori video}
I parametri di progetto che caratterizzano un codificatore video sono:
\begin{itemize}
	\item modalità di predizione temporale per la stima del movimento
	\item struttura del GOP e dipendenze tra fotogrammi
	\item strategia di mode selection per la predizione spaziale e temporale
\end{itemize}

\noindent
Gli indici di prestazione di un codificatore video sono:
\begin{itemize}
	\item bitrate: quantità di dati per unità di tempo necessaria per trasmettere il video
	\item qualità dell'immagine: misurata in termini di PSNR o tramite test soggettivi
	\item complessità computazionale: tempo di codifica e decodifica e risorse computazionali richieste
	\item ritardo di codifica e decodifica: tempo necessario per accumulare dati prima di poterli codificare
	o decodificare (es. per la predizione temporale rispetto ad immagini successive)
	\item robustezza agli errori: capacità di resistere a errori di trasmissione e perdita di pacchetti
\end{itemize}

\subsection{Stima e compensazione del movimento per block matching}
\subsubsection*{Suddivisione in blocchi}
Il movimento in un video può essere dovuto allo spostamento degli oggetti nella scena o al movimento della telecamera.
Per stimare correttamente il movimento dei singoli oggetti, l'immagine viene suddivisa in blocchi di dimensione \(N \times M\)
pixel e per ciascuno viene fatta la predizione del movimento. I blocchi si indicano attraverso la notazione \(B_{k}^{(\vec{p})}\)
in cui \(k\) è l'indice del fotogramma e \(\vec{p}(x,y)\) è la posizione del pixel centrale del blocco.

\subsubsection*{Stima del movimento con vettore e compensazione}
L'obiettivo della stima del movimento non è quello di tracciare il movimento degli oggetti, ma di trovare un blocco di
riferimento più simile a quello da predire per minimizzare il residuo di predizione.

La stima del movimento consiste nell'individuare il vettore di movimento \(\vec{v}(x,y)\) che collega il blocco da
predire \(B_k^{(\vec{p})}\) ad un altro blocco \(B_h^{(\vec{p} + \vec{v})}\) di un altro fotogramma \(h\) usato come
riferimento per la predizione tale da minimizzare il costo di predizione \(J(\vec{v})\).
\[J(\vec{v}) = d\left( B_k^{(\vec{p})}, B_h^{(\vec{p} + \vec{v})} \right) + \lambda \, R(\vec{v}) \qquad\qquad \vec{v}\,^* = \arg\min_{\vec{v} \in \mathcal{W}} J(\vec{v}\,^*)\]

La dissimilarità (o differenza) tra due blocchi \(d(\vec{v}) = d(B_k^{(\vec{p})}, B_h^{(\vec{p} + \vec{v})})\) in
funzione del vettore di movimento può essere calcolata come somma delle differenze assolute dei pixel (SAD) o come
somma dei quadrati delle differenze dei pixel (SSD).

Siccome il costo complessivo è dato sia dalla rappresentazione del residuo che da quella del vettore \(\vec{v}\),
la formula di costo \(J(\vec{v})\) include un termine di penalizzazione \(R(\vec{v})\) che rappresenta il costo
per rappresentare il vettore di movimento e un parametro \(\lambda\) che bilancia l'importanza dei due termini.

\subsubsection*{Algoritmi e pattern di ricerca del vettore di movimento}
L'algoritmo più semplice per la ricerca del vettore di movimento è la ricerca esaustiva (full search) che analizza il costo
di predizione per tutti i possibili blocchi nella finestra di ricerca \(\mathcal{W}\) e seleziona quello con il costo minimo.
Tale algoritmo è ottimale, ma molto complesso e costoso in termini di tempo di calcolo.

In alternativa esistono algoritmi euristici subottimali come la ricerca hexagon che espande la ricerca analizzando i blocchi
adiacenti e proseguendo solo nella direzione del blocco con il costo minore. La complessità computazionale è molto più bassa,
ma non viene garantita l'ottimalità della soluzione.

\subsubsection*{Codifica del residuo di predizione}
La codifica del residuo di predizione può avvenire tramite diverse codifiche in base al contesto:
\begin{itemize}
	\item codifica Exp-Golomb, efficace in generale siccome i residui hanno valori sparsi e vicini allo zero
	\item codifica aritmetica, più efficace se c'è poco movimento e quindi i residui sono molto vicini allo zero
	\item codifica JPEG semplificata, in generale siccome i residui hanno valori molto simili tra di loro e per molti blocchi
	è sufficiente solo il coefficiente DC
	\item codifica JPEG completa, se il residuo ha valori poco sparsificati ad esempio per cambi di scena
\end{itemize}

\subsubsection*{Parametri e metriche prestazione della ricerca block-matching}
I parametri di progetto che caratterizzano la stima del movimento sono:
\begin{itemize}
	\item forma e dimensione dei blocchi: blocchi piccoli riducono il costo dei residui di predizione, ma aumentano il
	numero di vettori di movimento da rappresentare e quindi anche costo complessivo
	\item dimensione della finestra di ricerca \(\mathcal{W}\): aumentando la finestra di ricerca aumenta la probabilità
	di trovare blocchi di riferimento più simili, ma aumenta anche la complessità computazionale
	\item funzione di dissimilarità: si preferisce la SAD in quanto più semplice da calcolare e simile al MSE
	\item strategia di ricerca: la ricerca esaustiva è la più accurata, ma anche la più complessa, per cui si preferiscono
	algoritmi euristici subottimali per ridurre la complessità computazionale
\end{itemize}
\noindent
Le prestazioni della stima del movimento si misurano in termini di:
\begin{itemize}
	\item accuratezza della stima: misurata in termini di PSNR o MSE del residuo di predizione
	\item bitrate o costo della predizione: misurato in termini di numero di bit necessari per rappresentare il residuo
	e il vettore di movimento
	\item complessità computazionale: misurata in termini di tempo di calcolo e risorse computazionali richieste
\end{itemize}

\subsubsection*{Criticità della stima del movimento}
La stima del movimento risulta poco efficace in prossimità dei bordi dell'immagine o in aree di disocclusione in cui
compaiono nuovi oggetti che non hanno corrispondenza nei fotogrammi precedenti.

Si nota inoltre che in aree senza strutture come il cielo o superfici uniformi, il block-matching si basa sui pattern di rumore
del blocco, che non sempre sono coerenti con il movimento reale. Tale problema, risulta irrilevante siccome il residuo di
predizione sarà comunque quasi nullo.

\subsection{Tipi di immagine e GOP}
\subsubsection*{Tipi di immagine/frame}
Le immagini o frame di un group of pictures si classificano in tre categorie principali:
\begin{itemize}
	\item frame I - intra coded: frame codificati senza predizione temporale, ma solo con predizione spaziale dalla stessa
	immagine o codifica JPEG per cui hanno bassa complessità di codifica/decodifica; sono usati come base per la predizione di
	tutti gli altri frame del GOP e garantiscono l'accesso casuale rapido al video e maggiore immunità agli errori
	\item frame P - predictive: frame codificati con predizione temporale da un frame precedente (I o P) oppure tramite
	predizione spaziale dalla stessa immagine, risultano di complessità maggiore rispetto ai frame I, ma hanno bitrate inferiore
	\item frame B - bidirectional predictive: frame codificati con predizione temporale da due frame, uno precedente e uno
	successivo (I o P), in questo modo si aumenta la probabilità di trovare blocchi di riferimento più simili a discapito
	di una maggiore complessità di codifica/decodifica e all'introduzione di latenza in quanto dipendono da frame successivi
\end{itemize}
I frame I e P sono detti anchor frame (AF) in quanto sono usati come riferimento per la predizione di altri frame (P o B).

\subsubsection*{Struttura di un Group-Of-Pictures}
Il GOP è una sequenza di immagini che inizia con un frame I e termina con il frame precedente al successivo frame I.
Un GOP è definito da due parametri principali:
\begin{itemize}
	\item \(N\): intervallo da un frame I al successivo frame I, denota il numero di frame nel GOP
	\item \(M\): intervallo tra due AF successivi, denota il numero di frame B tra due I o P successivi
\end{itemize}

\subsubsection*{Tasso e distorsione dei vari tipi di immagine}
Si osserva che i fotogrammi I hanno un bitrate di 3-5 volte maggiore rispetto ai fotogrammi P e 10-20 volte maggiore rispetto
ai fotogrammi B. La qualità dei fotogrammi I è forzatamente maggiore rispetto ai fotogrammi P e B in quanto sono usati come
riferimento per l'intero GOP.

I parametri \(N\) e \(M\) del GOP fungono da compromesso tra il bitrate complessivo e la complessità di codifica/decodifica.
Avere tanti AF aumenta il bitrate complessivo, ma riduce la complessità di codifica/decodifica, viceversa aumentando il numero
di frame B si riduce il bitrate, ma si aumenta la complessità di codifica/decodifica e la latenza.

\subsection{Codificatori e decodificatori ibridi}
\subsubsection*{Codificatore ibrido per frame intra}
Il codificatore ibrido per frame intra è composto da:
\begin{itemize}
	\item catena di blocchi per la codifica simil-JPEG (DCT, quantizzazione Q e variable length coding VLC)
	\item anello di loopback che ricostruisce il frame decodificato tramite la normalizzazione dei coefficienti quantizzati
	(Q*) e la Inverse DCT tipici del decodificatore JPEG ed effettua una predizione spaziale intra-frame tra i blocchi del
	frame corrente per ridurre la ridondanza spaziale tra blocchi adiacenti
\end{itemize}

\begin{figure}[htbp]
	\centering
	\begin{tikzpicture}[
		% Stile per i blocchi di elaborazione generici
		block/.style={
			rectangle,
			draw=blue!70!black,
			fill=blue!5,
			thick,
			align=center,
			minimum width=1.1cm,
			rounded corners=3pt,
			font=\footnotesize
		},
		% Stile specifico per il canale di trasmissione
		channel/.style={
			rectangle,
			draw=orange!85!black,
			fill=orange!8,
			thick,
			align=center,
			rounded corners=3pt,
			font=\footnotesize
		},
		% Stile per i sommatori
		sum/.style={
			circle,
			draw=blue!70!black,
			fill=blue!5,
			thick,
			minimum size=0.4cm,
			inner sep=0pt
		},
		% Stile per le frecce e le linee
		arrow/.style={
			-{Stealth[scale=1.2]},
			thick
		},
		line/.style={
			thick
		},
		% Stile per i segnali testuali
		signal/.style={
			align=center,
			font=\footnotesize
		}
	]

		% Catena di blocchi
		\node (in) at (-0.5,0) [signal] {\textbf{Input}};
		\node[sum] (adder1) at (1.5,0) {}; 	% Sommatore 1
		\draw[thick, blue!70!black] (adder1.north) -- (adder1.south);
		\draw[thick, blue!70!black] (adder1.west) -- (adder1.east);
		\node[block] (dct) at (3,0) {DCT};
		\node[block] (q) at (5,0) {Q};
		\coordinate (split1) at (6.3,0);
		\fill[blue!70!black] (split1) circle (2pt); % Punto di derivazione (dot)
		\node[block] (vlc) at (7.5,0) {VLC};
		\node[block] (cbuf) at (9.5, 0) {Channel\\Buffer};
		\node (out) at (12, 0) [signal] {\textbf{Output}};

		\node[block] (iq) at (6.3, -1) {Q*};
		\node[block] (idct) at (6.3, -2) {IDCT};
		\node[sum] (adder2) at (6.3, -3) {}; % Sommatore 2
		\draw[thick, blue!70!black] (adder2.north) -- (adder2.south);
		\draw[thick, blue!70!black] (adder2.west) -- (adder2.east);
		\node[block] (framebuf) at (3.3, -3) {Decoded\\Frame Buffer};
		\node[block] (intrapred) at (3.3, -1.5) {Intra\\Prediction};

		% Collegamenti ed etichette
		\draw[arrow] (in) -- (adder1.west);
		\node[signal] at ($(adder1.west)+(-0.7, 0.2)$) {$B_2$};
		\node[signal] at ($(adder1.west)+(-0.2, 0.25)$) {$+$};
		\node[signal] at ($(adder1.west)+(0.8, -0.5)$) {$\hat{B}_2$};
		\node[signal] at ($(adder1.south)+(-0.25, -0.2)$) {$-$};
		\node[signal] at ($(adder1.east)+(0.3, 0.2)$) {$e_2$};

		\draw[arrow] (adder1.east) -- (dct.west);
		\draw[arrow] (dct.east) -- (q.west);
		\draw[line] (q.east) -- (split1);
		\draw[arrow] (split1) -- (vlc.west);
		\draw[arrow] (vlc.east) -- (cbuf.west);
		\draw[arrow] (cbuf.east) -- (out);

		\draw[arrow] (split1) -- (iq.north);
		\draw[arrow] (iq.south) -- (idct.north);
		\draw[arrow] (idct.south) -- (adder2.north);
		\node[signal] at ($(adder2.north)+(0.3, 0.3)$) {$\hat{e}_2$};
		\node[signal] at ($(adder2.west)+(-0.4, 0.2)$) {$\hat{B}_2$};

		\draw[arrow] (adder2.west) -- (framebuf.east);
		\draw[arrow] (framebuf.north) -- (intrapred.south);
		\coordinate (split2) at ($(intrapred.north -| adder1.south)+(0, 0.4)$);
		\draw[line] (intrapred.north) -- +(0, 0.4) -- (split2);
		\draw[arrow] (split2) -- (adder1.south);
		\draw[arrow] (split2) -- +(0,-3) -| (adder2.south);
	\end{tikzpicture}
\end{figure}

\subsubsection*{Codificatore ibrido per frame inter}
Il codificatore ibrido per frame inter è composto da:
\begin{itemize}
	\item la stessa base del codificatore di tipo intra con codifica simil-JPEG e predizione spaziale intra-frame
	\item un blocco di motion estimation (ME) che calcola il vettore di movimento \(\vec{v}\)
	\item un blocco di motion compensation (MC) che calcola il blocco predetto \(\hat{B}_k^{(\vec{p})}\) a partire dal
	vecchio frame di riferimento nel decoded frame buffer e dal vettore di movimento \(\vec{v}\)
	\item un blocco di selezione tra predizione spaziale di tipo intra e predizione temporale di tipo inter
\end{itemize}

\noindent
La scelta tra l'utilizzo della predizione spaziale o quella temporale dipende da un criterio di valutazione del
compromesso tra bitrate e distorsione come il seguente:
\[m_\text{modo di codifica} \in \left\{ \text{INTRA}, \text{INTER} \right\} \qquad\qquad m^* = \arg\min_{m} \! \underbrace{\, D_m \,}_\text{distorsione} \! + \; \lambda \! \underbrace{\, R_m \,}_\text{bitrate} \]
Lo standard non definisce il tipo di criterio da utilizzare, la scelta che viene lasciata al progettista del
codificatore che potrebbe anche pensare di usare sempre solo una delle due predizioni per risparmiare tempo
di calcolo. In questo modo si permette di avere concorrenza tra diversi produttori.


\begin{figure}[htbp]
	\centering
	\begin{tikzpicture}[
		% Stile per i blocchi di elaborazione generici
		block/.style={
			rectangle,
			draw=blue!70!black,
			fill=blue!5,
			thick,
			align=center,
			minimum width=1.1cm,
			rounded corners=3pt,
			font=\footnotesize
		},
		% Stile specifico per il canale di trasmissione
		channel/.style={
			rectangle,
			draw=orange!85!black,
			fill=orange!8,
			thick,
			align=center,
			rounded corners=3pt,
			font=\footnotesize
		},
		% Stile per i sommatori
		sum/.style={
			circle,
			draw=blue!70!black,
			fill=blue!5,
			thick,
			minimum size=0.4cm,
			inner sep=0pt
		},
		% Stile per le frecce e le linee
		arrow/.style={
			-{Stealth[scale=1.2]},
			thick
		},
		line/.style={
			thick
		},
		% Stile per i segnali testuali
		signal/.style={
			align=center,
			font=\footnotesize
		}
	]
		% Catena di blocchi
		\node (in) at (0,0) [signal] {\textbf{Input}};
		\node[sum] (adder1) at (2,0) {}; 	% Sommatore 1
		\draw[thick, blue!70!black] (adder1.north) -- (adder1.south);
		\draw[thick, blue!70!black] (adder1.west) -- (adder1.east);
		\node[block] (dct) at (4,0) {DCT};
		\node[block] (q) at (6.5,0) {Q};
		\coordinate (split1) at (8,0);
		\fill[blue!70!black] (split1) circle (2pt); % Punto di derivazione (dot)
		\node[block] (vlc) at (9.5,0) {VLC};
		\node[block] (cbuf) at (11.5,0) {Channel\\Buffer};
		\node (out) at (13.5, 0) [signal] {\textbf{Output}};

		\node[block] (iq) at (8, -1) {Q*};
		\node[block] (idct) at (8, -2) {IDCT};
		\node[sum] (adder2) at (8, -3) {}; % Sommatore 2
		\draw[thick, blue!70!black] (adder2.north) -- (adder2.south);
		\draw[thick, blue!70!black] (adder2.west) -- (adder2.east);
		\node[block] (framebuf) at (6, -3) {Decoded\\Frame Buffer};
		\node[block] (intrapred) at (6, -1.2) {Intra\\Prediction};
		\node[block] (interintra) at (3.5, -1.2) {Inter/\\Intra};
		\node[block] (mc) at (3.5, -3) {Motion\\Compensation};

		\node (in2) at (-0.3,-3.3) [signal] {\textbf{old frame}\\(reference)\\dal DFB};
		\node[block] (me) at (1, -2) {Motion\\Estimation};

		% Collegamenti ed etichette
		\draw[arrow] (in) -- (adder1.west);
		\node[signal] at ($(adder1.west)+(-0.7, 0.2)$) {$B_2$};
		\node[signal] at ($(adder1.west)+(-0.2, 0.25)$) {$+$};
		\node[signal] at ($(adder1.south)+(0.25, -0.5)$) {$\hat{B}_2$};
		\node[signal] at ($(adder1.south)+(-0.25, -0.2)$) {$-$};
		\node[signal] at ($(adder1.east)+(0.3, 0.2)$) {$e_2$};

		\draw[arrow] (adder1.east) -- (dct.west);
		\draw[arrow] (dct.east) -- (q.west);
		\draw[line] (q.east) -- (split1);
		\draw[arrow] (split1) -- (vlc.west);
		\draw[arrow] (vlc.east) -- (cbuf.west);
		\draw[arrow] (cbuf.east) -- (out);

		\draw[arrow] (split1) -- (iq.north);
		\draw[arrow] (iq.south) -- (idct.north);
		\draw[arrow] (idct.south) -- (adder2.north);
		\node[signal] at ($(adder2.north)+(0.3, 0.3)$) {$\hat{e}_2$};
		\node[signal] at ($(adder2.west)+(-0.3, 0.2)$) {$\hat{B}_2$};

		\draw[arrow] (adder2.west) -- (framebuf.east);
		\draw[arrow] (framebuf.north) -- (intrapred.south);
		\draw[arrow] (intrapred.west) -- (interintra.east);
		\draw[arrow] (framebuf.west) -- (mc.east);
		\draw[arrow] (mc.north) -- (interintra.south);
		\coordinate (split2) at ($(interintra.west -| adder1.south)$);
		\draw[line] (interintra.west) -- (split2);
		\draw[arrow] (split2) -- (adder1.south);
		\draw[arrow] (split2) -- +(0,-2.6) -- ($(split2 -| adder2.south)+(0,-2.6)$) to[arrow] (adder2.south);

		\draw[arrow] (1,0) -- (me.north);
		\coordinate (split3) at ($(me.south -| mc.south)+(0,-1.7)$);
		\draw[line] (me.south) -- +(0,-1.7) -- node[signal, midway, shift={(0cm, 0.03cm)}, fill=white, inner sep=1pt] {motion vector} (split3) {};
		\draw[arrow] (split3.south) -- (mc.south);
		\draw[arrow] (split3.south) -| (vlc.south);

		\draw[arrow] (in2.north) |- (me.west);
	\end{tikzpicture}
\end{figure}

\subsubsection*{Decodificatore ibrido}
Il decodificatore ibrido ripercorre al contrario la catena di elaborazione del codificatore ibrido. Risulterà più
veloce rispetto ai codificatori quanto non deve effettuare calcoli sulla predizione del movimento o selezioni tra
predizione spaziale e temporale.

\begin{figure}[htbp]
	\centering
	\begin{tikzpicture}[
		% Stile per i blocchi di elaborazione generici
		block/.style={
			rectangle,
			draw=blue!70!black,
			fill=blue!5,
			thick,
			align=center,
			minimum width=1.1cm,
			rounded corners=3pt,
			font=\footnotesize
		},
		% Stile specifico per il canale di trasmissione
		channel/.style={
			rectangle,
			draw=orange!85!black,
			fill=orange!8,
			thick,
			align=center,
			rounded corners=3pt,
			font=\footnotesize
		},
		% Stile per i sommatori
		sum/.style={
			circle,
			draw=blue!70!black,
			fill=blue!5,
			thick,
			minimum size=0.4cm,
			inner sep=0pt
		},
		% Stile per le frecce e le linee
		arrow/.style={
			-{Stealth[scale=1.2]},
			thick
		},
		line/.style={
			thick
		},
		% Stile per i segnali testuali
		signal/.style={
			align=center,
			font=\footnotesize
		}
	]
		% Catena di blocchi
		\node (in) at (0,0) [signal] {\textbf{Input}};
		\node[block] (cbuf) at (2.5,0) {Channel\\buffer};
		\node[block] (vld) at (5,0) {VLD};
		\node[block] (iq) at (7.5,0) {Q*};
		\node[block] (idct) at (10,0) {IDCT};
		\node[sum] (adder) at (10,-1.5) {};
		\draw[thick, blue!70!black] (adder.north) -- (adder.south);
		\draw[thick, blue!70!black] (adder.west) -- (adder.east);
		\node (out) at (12,-1.5) [signal] {\textbf{Output}};
		\node[block] (framebuf) at (7.5,-1.5) {Decoded\\Frame Buffer};
		\node[block] (intrapred) at (7.5,-3) {Intra\\Prediction};
		\node[block] (switch) at (10,-3) {Intra/Inter};
		\node[block] (mc) at (3.5,-3) {Motion\\Compensation};

		% Collegamenti ed etichette
		\draw[arrow] (in) -- (cbuf.west);
		\draw[arrow] (cbuf.east) -- (vld.west);
		\draw[arrow] (vld.east) -- (iq.west);
		\draw[arrow] (iq.east) -- (idct.west);
		\draw[arrow] (idct.south) -- (adder.north);
		\draw[arrow] (adder.east) -- (out.west);
		\node[signal] at ($(adder.north)+(0.15,0.5)$) {$\hat{e}$};
		\node[signal] at ($(adder.east)+(0.5,0.2)$) {$\hat{B}$};
		\node[signal] at ($(adder.west)+(-0.5,0.2)$) {$\hat{B}$};
		\node[signal] at ($(adder.south)+(0.15,-0.5)$) {$\tilde{B}$};
		\draw[arrow] (adder.west) -- (framebuf.east);
		\draw[arrow] (framebuf.west) -| (mc.north);
		\draw[arrow] (mc.south) -- ++(0,-0.8) -| ($(switch.south)+(0.15,0)$);
		\draw[arrow] (switch.north) -- (adder.south);
		\draw[arrow] (framebuf.south) -- (intrapred.north);
		\draw[arrow] (intrapred.east) -- (switch.west);
		\draw[arrow] ($(vld.south)+(-0.1,0)$) -- ++(0,-0.9) -- node[signal, midway, shift={(0cm, 0.03cm)}, fill=white, inner sep=1pt] {motion vector} ++(-3,0) |- (mc.west);
		\draw[arrow] ($(vld.south)+(0.1,0)$) -- ++(0,-3.7) -| node[signal, midway, shift={(-2.2cm, 0.01cm)}, fill=white, inner sep=1pt] {inter/intra parameter} ($(switch.south)+(-0.15,0)$);
	\end{tikzpicture}
\end{figure}

\newpage

\subsection{Standard di codifica video principali}
\subsubsection*{Evoluzione degli standard di codifica video}
\begin{itemize}[itemsep=2pt]
	\item \textbf{1995 - MPEG-2}: \\
	primo standard video di grande diffusione usato per DVD e TV satellitare, è composto da un codec ibrido (I,P,B) senza
	predizione spaziale (è presente solo JPEG), supporta fino a 1080p a 60fps con bitrate di 80 Mbps, fornisce strumenti di
	codifica scalabile ovvero è possibile estrarre più flussi video con qualità ridotte da un unico flusso codificato (utile
	per adattare la trasmissione alla rete)
	\item \textbf{2003 - H.264/AVC}: \\
	standard ancora oggi più diffuso e usato per lo streaming, rispetto a MPEG-2 permette predizione spaziale, blocchi a
	dimensione variabile, trasformate simil-DCT intere e viene aggiunto un deblocking filter a fine codifica per rimuovere
	gli artefatti di quantizzazione, il bitrate viene dimezzato rispetto a MPEG-2
	\item \textbf{2013 - H.265/HEVC}: \\
	usato per streaming 4K e DVB-T2, è composto da un codec ibrido basato su alberi quaternari, il bitrate viene dimezzato
	rispetto a H.264, non è molto diffuso perché protetto da licenze e royalties
	\item \textbf{2018 - AV1}: \\
	alternativa royalty-free al H.265, usato per streaming Youtube e Netflix
	\item \textbf{2021 - VVC/H.266}: \\
	usato per video 8K e VR, è composto da un versatile video coding che aumenta la alta complessità di \(10\times\) rispetto
	a H.265, dimezzandone il bitrate
	\item \textbf{2025+ - Neural Video Coding}: \\
	nuovi codec basati su reti neurali (DNN e autoencoders) contenuti nei nuovi standard MPEG-AI
\end{itemize}

\subsubsection*{Digitale terrestre}
\begin{itemize}
	\item \textbf{H.265/HEVC} 55-60\% in crescita: standard per il DVB-T2 in Europa (diventato standard nell'ottobre 2025)
	e ATSC 3.0 negli USA, necessario per le trasmissioni in 4K e HDR
	\item \textbf{H.264/AVC} 30-35\% in calo: usato come fallback per i canali SD/HD o in televisori più vecchi
	\item \textbf{MPEG-2} 5-10\% in calo: limitato a mercati di nicchia o canali locali
\end{itemize}

\subsubsection*{Streaming}
\begin{itemize}
	\item \textbf{AV1} 35-40\%: adottato massicciamente da Netflix, Youtube e Meta
	\item \textbf{H.264/AVC} 30\%: utilizzato per retrocompatibilità con dispositivi più vecchi
	\item \textbf{H.265/HEVC} 25\%: per contenuti 4K HDR su dispositivi Apple e TV di fascia alta
\end{itemize}

\subsubsection*{Video real-time}
\begin{itemize}
	\item \textbf{H.264/AVC} 50\%: molto usato grazie all'accelerazione hardware molto diffusa
	\item \textbf{VP9} 25\%: versione royalty free di H.264, utilizzato da Google Meet e in applicazioni browser-based grazie
	alla sua efficienza superiore rispetto a H.264
	\item \textbf{VP8} 15\%: ancora presente in WebRTC per la bassissima complessità computazionale
	\item \textbf{AV1} 10\% in crescita: utilizzato per scenari di banda estremamente limitata, grazie al basso bitrate
\end{itemize}
